\chapter{Introducción}
\label{sec:arq-introduccion}

GraphHunter es un motor de caza de amenazas sobre grafos temporales
construidos a partir de logs de seguridad. Este libro describe
\emph{cómo está construido}: qué capas lo componen, qué piezas viven
en cada capa, por qué se eligieron esos límites y qué patrones se
repiten en todo el código.

\textbf{Véase también} el libro \emph{Diseño completo}
(\filepath{docs/diseno/diseno.pdf}) para una guía pedagógica
del diseño algorítmico: pipeline de ingesta, matcher tier-por-tier
(DFS / LFTJ-K4 / Yannakakis chain / Yannakakis tree-walk), índices
auxiliares (NLF, k-hop reachability), DSL, y el plano de obra de
las stages futuras L3 stage 2b/2c.

\section{Audiencia}

Está pensado para tres lectores distintos:

\begin{itemize}
    \item \textbf{Nuevo contribuidor}: necesita el mapa del
    repositorio y el razonamiento detrás de la estructura antes de
    abrir su primer pull request.
    \item \textbf{Revisor arquitectónico}: evalúa propuestas de
    cambio que tocan múltiples capas, o sanciona nuevas dependencias
    entre crates.
    \item \textbf{Operador técnico}: opera el sistema en producción
    y necesita entender los límites de aislamiento (FFI, canonical
    API, model gate) para diagnosticar problemas.
\end{itemize}

\section{Alcance}

Este libro cubre \emph{la arquitectura en sí}: estructura,
decisiones, patrones, deudas. No cubre:

\begin{itemize}
    \item Las \textbf{superficies externas} (HTTP API, MCP tools,
    DSL, C ABI) — eso vive en el libro de Especificación.
    \item El \textbf{flujo de uso} del sistema (ingestar, cazar,
    anotar, exportar) — eso vive en el Manual.
    \item La \textbf{justificación experimental} de algoritmos
    (matcher SIMD, DFS temporal, scoring compuesto) — eso vive en el
    Paper.
\end{itemize}

\section{Mapa del libro}

\begin{itemize}
    \item El capítulo~\ref{sec:shearing-layers} introduce el modelo
    de capas de cizalla (\emph{shearing layers}) que organiza el
    proyecto.
    \item El capítulo~\ref{sec:grafo-deps} muestra el grafo de
    dependencias inter-crate y las propiedades que garantiza.
    \item Los capítulos~\ref{sec:capa-core}, \ref{sec:capa-platform}
    y~\ref{sec:capa-apps} recorren cada capa en detalle.
    \item El capítulo~\ref{sec:patrones} describe patrones
    transversales: typestate, RAII, canonical API, model gate.
    \item El capítulo~\ref{sec:decisiones} resume las seis decisiones
    arquitectónicas (ADR-001 a ADR-006) que fundamentan el estado
    actual.
    \item El capítulo~\ref{sec:deudas} documenta honestamente la
    deuda conocida y el trabajo futuro.
    \item El apéndice~\ref{sec:adrs} contiene las ADRs textualmente.
    El apéndice~\ref{sec:metricas} reporta métricas cuantitativas
    (LOC, tests, crates, ciclos).
\end{itemize}

\section{Convenciones}

\begin{itemize}
    \item \filepath{rutas/de/archivo} en tipografía monoespaciada.
    \item \code{identificadores_de_código} igualmente monoespaciados.
    \item Crates Rust se escriben con su path en el repositorio:
    \filepath{core/graph-engine} significa ``el crate que vive en
    \filepath{core/graph-engine/Cargo.toml}''.
    \item Cuando un párrafo dice ``ADR-00N'' se refiere al apéndice
    \ref{sec:adrs}, que contiene el documento completo.
\end{itemize}
